\section{Revisão da Literatura}
\label{sec:literatura}

% ESCREVER PARÁGRAFO DE ABERTURA.

% Origem: introdução, REF1--REF2 e fundamentação conceitual de REF7/REF11.
\subsection{Expansão agropecuária, retornos privados, custos públicos e os casos da Amazônia, Cerrado e Pantanal}

A expansão agropecuária envolve escolha econômica entre usos alternativos da terra, na qual a conversão da vegetação em área produtiva para lavoura ou pecuária pode ampliar a renda privada associada à produção, ao mesmo tempo em que reduz estoques de capital natural e serviços ecossistêmicos que nem sempre são incorporados aos retornos percebidos pelos agentes. Essa diferença entre o valor social da vegetação e os benefícios que o proprietário consegue apropriar integra o enquadramento de \citeonline{Angelsen2010}, que distingue as rendas da agricultura e da floresta. A decisão depende, portanto, da comparação entre alternativas efetivamente acessíveis ao agente, e não apenas do valor da produção que poderá ser obtida na área convertida.

Nesse sentido, os efeitos da conversão podem ultrapassar os limites da propriedade convertida, embora a extensão desses efeitos dependa dos processos considerados. No Cerrado brasileiro, por exemplo, \citeonline{Schwamback2024} documentaram maior escoamento de nutrientes sob pastagem e maior perda de solo sob cana, em comparação com o cerrado arborizado. Na fronteira Amazônia-Cerrado, a substituição de floresta por soja esteve associada à redução do carbono do solo em até 100 cm de profundidade, com respostas diferentes conforme a profundidade e o manejo \cite{Popin2025}. Esses estudos documentam alterações biofísicas, mas não mensuram, por si sós, o custo social total da conversão.

Além das alterações biofísicas documentadas, a conversão pode ter implicações econômicas ao modificar serviços do solo que sustentam a produção ao longo do tempo. Nessa perspectiva, \citeonline{Brady2015} propuseram valorar esses serviços a partir de seus efeitos sobre fluxos futuros de renda agrícola, utilizando dados de experimentos de longo prazo na Europa e nos Estados Unidos que comparam a produtividade agrícola sob diferentes práticas de adubação e manejo do solo.

Com efeito, o problema permanece ainda menos resolvido quando se comparam benefícios e custos privados, como o lucro de um fazendeiro, e públicos, como os impactos gerados para outros agentes e territórios. Isso ocorre porque os serviços ambientais diferem entre lugares, operam em escalas distintas e muitos não possuem preços diretamente observáveis. Estudos que integram informações econômicas e biofísicas mostram a importância da localização e das alternativas de uso da terra para valorar esses serviços e orientar a conservação \cite{Crossman2009,Gascoigne2011}. Assim, a oposição entre conservar e produzir pode refletir diferenças entre retornos apropriados privadamente e custos distribuídos socialmente.

Para a Amazônia brasileira, \citeonline{Mann2012} estimaram esse descompasso ao comparar o retorno econômico potencial da conversão da floresta para a produção de soja com o valor monetário dos serviços ecossistêmicos mantidos pela vegetação nativa. Para tanto, combinaram estimativas do retorno da soja calculadas por um modelo de produtividade e pela diferença entre receitas e custos, com a valoração dos serviços ecossistêmicos por preços-sombra derivados de uma função de produção Cobb-Douglas, que incorporou a produtividade primária da vegetação como indicador desses serviços. Sob as hipóteses adotadas, a conclusão foi que o valor dos serviços supera a renda da soja em parcela expressiva da área analisada, e a omissão desses serviços reduz o bem-estar social relativamente ao ótimo modelado.

No caso do Pantanal, existem tentativas de mensurar economicamente os serviços ambientais, mas os próprios resultados evidenciam dificuldades de escala, cobertura e comparabilidade. \citeonline[p.~3]{SeidlMoraes2000}, ao recalcularem para a sub-região da Nhecolândia \footnote{Região do Pantanal sul-mato-grossense (Bacia do Alto Paraguai) caracterizada por florestas sazonais, pastagens nativas e lagoas alcalinas, com tradicional manejo da pecuária bovina extensiva.} estimativas globais com dados locais, obtiveram valor anual de aproximadamente US\$ 5,8 mil por hectare, em dólares de 1994, equivalente a cerca de metade daquele derivado dos parâmetros globais, diferença atribuída às condições biofísicas regionais. 

Mais recentemente, \citeonline[p.~7]{FairbrassEtAl2025}, ao estruturarem indicadores de capital natural para o Pantanal, buscaram reunir informações sobre a extensão, o estado de conservação e o valor econômico dos ecossistemas, os serviços ambientais oferecidos e as atividades humanas que deles se beneficiam. Dos 29 indicadores previstos para os ativos ecossistêmicos, encontraram dados que atendiam aos critérios do estudo para apenas 14, menos da metade do conjunto. Essa cobertura incompleta limitou a caracterização das relações entre a economia e o ambiente, dado que, embora houvesse informações sobre atividades com valor de mercado, como pecuária, soja, pesca e turismo, faltavam indicadores sobre a condição e o valor econômico dos ecossistemas que as sustentam, sobre parte das emissões geradas por essas atividades e acerca dos processos ecológicos essenciais, como a ciclagem de nutrientes e a manutenção da biodiversidade.

Sublinhe-se que as dificuldades de cobertura observadas no caso do Pantanal no parágrafo anterior também se estendem a outros biomas. \citeonline{CampoliStivali2026} reuniram estudos de valoração econômica da vegetação nativa nos seis biomas brasileiros (Amazônia, Cerrado, Pantanal, Mata Atlântica, Caatinga e Pampa) e padronizaram estimativas de custo social do desmatamento em dólares de 2025 por hectare. Ao final, os autores concluíram que os estudos diferem tanto em relação aos benefícios da vegetação incluídos na valoração econômica quanto nos métodos empregados. 

Como se percebe, as diferenças entre os valores estimados para os ganhos econômicos ambientais podem refletir tanto as características dos biomas como as diferenças no que cada estudo conseguiu mensurar. Em conjunto, os trabalhos discutidos mostram por que o retorno privado da conversão não representa, necessariamente, seu resultado econômico para a sociedade: a renda obtida com o novo uso da terra pode coexistir com perdas de serviços ambientais cujos custos não recaem integralmente sobre quem decide converter. Essa distinção fundamenta o recorte deste artigo, voltado a investigar como os incentivos privados se relacionam com a conversão da vegetação nativa. Cabe, portanto, examinar como os preços agropecuários e a capacidade produtiva da terra modificam a atratividade dos usos alternativos disponíveis aos produtores.

Nessa perspectiva, a contribuição deste artigo situa-se na análise das decisões que transformam o estoque ambiental, e não na estimação do valor monetário de seus serviços. Busca-se esclarecer como incentivos privados e condições institucionais se relacionam com a conversão para diferentes usos agropecuários. Essa delimitação aproxima a discussão sobre o descompasso entre retornos privados e custos sociais de uma questão empiricamente observável, sem supor que as lacunas de valoração serão resolvidas pelo exame de H1.

% Origem: REF3, REF6--REF8 e diálogo empírico de REF11.
\subsection{Preços, capacidade produtiva e usos alternativos da terra}

No argumento deste artigo, os preços e a capacidade produtiva representam dimensões observáveis dos incentivos ao uso agropecuário. Mantidos os custos e as demais condições, uma elevação de preço representa maior acréscimo de receita por hectare onde a produção potencial do uso beneficiado é maior, sendo isso justamente a relação econômica que motiva H1. Sua tradução em conversão efetiva, contudo, depende da vegetação disponível, dos custos de abertura da área, das alternativas em terras já convertidas e das restrições institucionais. Trata-se de uma previsão condicional a investigar, não de uma consequência necessária de qualquer aumento de produtividade.

Conforme suscitado logo no início desta seção, a conversão da vegetação nativa pode ser analisada como resultado de escolhas econômicas condicionadas por retornos relativos. \citeonline{Angelsen2010} organiza as políticas de redução do desmatamento em torno das rendas relativas da agricultura e da floresta, distinguindo intervenções que reduzem a atratividade da expansão, elevam os benefícios apropriáveis da conservação ou restringem diretamente a conversão. Esse enquadramento permite compreender por que a intensificação produtiva pode ampliar a pressão local de expansão em determinadas condições, sem presumir que produzir mais implique necessariamente converter mais terra.

A evidência de \citeonline{LubowskiEtAl2008}, para os Estados Unidos, relaciona os retornos líquidos dos usos alternativos e a qualidade da terra às decisões dos proprietários. Na Amazônia brasileira, \citeonline{Pfaff1999} identifica a relevância das características da terra e dos custos de transporte, advertindo que os efeitos estimados de políticas podem refletir a própria seleção de onde elas são implementadas. \citeonline{ChomitzThomas2003}, por sua vez, analisam conversão e intensidade pecuária em escala espacial detalhada, considerando precipitação, solos, infraestrutura, acesso a mercados e proteção territorial. Essas contribuições sustentam a consideração das condições produtivas e de acesso na análise de uso da terra. A heterogeneidade espacial que documentam não equivale, por si só, a um teste da resposta diferencial a uma mesma variação de preços.

A distinção entre capacidade produtiva e produtividade observada é pertinente a essa passagem. O rendimento obtido em áreas já cultivadas também reflete decisões de manejo e utilização de insumos, além das características da terra selecionada para produzir. Uma medida de capacidade produtiva predeterminada precisa, portanto, ser avaliada quanto ao que representa e à sua posição temporal em relação à conversão. Da mesma forma, preços e produtividade, isoladamente, não revelam o benefício privado líquido esperado, pois não incorporam todos os custos, restrições e expectativas que integram a decisão do produtor.

Em escala tropical, \citeonline{BermanEtAl2023} combinaram preços internacionais e aptidão agrícola para examinar a relação entre incentivos de mercado e desmatamento entre 2001 e 2018.\footnote{A pesquisa abrangeu regiões tropicais de diferentes países, com destaque para a Amazônia, o Sudeste Asiático e a Bacia do Congo, onde a associação entre preços agrícolas e desmatamento foi mais forte.} Os autores encontraram uma associação positiva entre preços agrícolas e desmatamento, cuja intensidade variou conforme as características econômicas locais. Essa abordagem é pertinente à H1 porque permite considerar que um mesmo aumento de preço pode gerar incentivos distintos à conversão, dependendo das possibilidades produtivas de cada território. O uso da aptidão agrícola para medir a exposição aos preços, contudo, não demonstra, por si só, que a capacidade produtiva local modifica a resposta da conversão da forma prevista nesta pesquisa. Há também um desafio de identificação: os preços internacionais podem refletir choques de oferta, e é necessário justificar por que a exposição a suas variações pode ser considerada exógena ao desmatamento analisado. Portanto, o uso de preços internacionais não basta para assegurar uma interpretação causal dos resultados.

Sem embargo, a relação entre preços e conversão depende do uso que se expande e daquele que é substituído. Na bacia do Tapajós, \citeonline{Braganca2018} encontrou que a elevação do preço relativo de lavouras em relação à carne favoreceu a substituição de pastagens por cultivos (\textit{pasture-to-cropland}) e reduziu o desmatamento local. Uma mudança nos preços relativos entre atividades, portanto, não se confunde com uma elevação do retorno esperado da conversão da vegetação nativa. A pressão para abrir novas áreas agropastoris deve, assim, ser distinguida da reorganização da produção em áreas já convertidas. Para este artigo, essas margens de ajuste devem também ser distinguidas da intensificação produtiva. Soja, milho e pecuária tampouco podem ser combinados em um único índice sem explicitar unidades, exposição e mecanismo econômico.

À luz desses antecedentes, a contribuição aqui pretendida consiste em examinar as condições territoriais em que os incentivos de preços se associam à conversão da vegetação nativa, avançando sobre a relação documentada por \citeonline{BermanEtAl2023}, \citeonline{AssuncaoEtAl2015} e \citeonline{HardingEtAl2021}, dentre outros. O diferencial do artigo está em analisar conjuntamente a capacidade produtiva local, o ambiente institucional e os destinos da conversão, distinguindo a abertura de áreas para agricultura e pastagem da reorganização de usos em áreas anteriormente convertidas, destacada por \citeonline{Braganca2018}. Essa abordagem busca esclarecer por que movimentos semelhantes de preços podem se traduzir em respostas distintas de conversão e, com isso, oferecer evidências sobre o alcance e os limites das explicações existentes. Para o caso do estado de Mato Grosso, a contribuição pretendida é identificar quais atividades e territórios apresentam maior sensibilidade aos preços e em que medida essa resposta varia com as condições produtivas e institucionais, oferecendo subsídios à discussão de políticas de conservação que considerem essas diferenças, o que permitirá avaliar a adequação territorial de instrumentos regulatórios e econômicos de conservação, inclusive em outros contextos, regiões e biomas.

% Origem: REF1--REF2 e REF9; limites de escopo estabelecidos ao final de REF11.
\subsection{Instituições e respostas aos incentivos}

No caso do trade-off entre conversão indiscriminada e uso limitado do meio ambiente, as instituições pertinentes a um ambiente produtivo ajudam a compreender por que os produtores podem responder de formas distintas aos incentivos de preços discutidos anteriormente. Ao definir direitos de uso da terra, restrições à conversão e consequências de seu cumprimento ou descumprimento, as regras modificam os custos, os riscos e os retornos esperados das alternativas disponíveis. Sua influência sobre as decisões produtivas depende, porém, de como são implementadas em cada território. Examinar as respostas aos incentivos exige, portanto, considerar tanto o conteúdo das regras quanto os meios para aplicá-las e as possibilidades de ajuste dos produtores, seja pela abertura de novas áreas, seja pela reorganização de usos em terras já convertidas.

Para H1, importa distinguir a existência de uma regra, os meios disponíveis para implementá-la e as consequências que os agentes esperam de seu cumprimento ou descumprimento. Monitoramento, fiscalização e sanção correspondem a etapas relacionadas, mas distintas. O mecanismo proposto é que a possibilidade de detecção e a aplicação de consequências alterem o custo esperado da conversão irregular e, por essa via, a resposta aos incentivos econômicos. A evidência de fiscalização discutida a seguir oferece sustentação a esse mecanismo em contextos específicos, mas não autoriza atribuir o mesmo efeito a toda instituição relacionada ao uso da terra.

Sob essa perspectiva, o enquadramento econômico não supõe efeitos uniformes das instituições. \citeonline{PachecoMeyer2022} comparam regimes fundiários em diferentes contextos brasileiros e mostram que os resultados dependem do regime de origem, da alternativa de destinação da terra e do contexto socioambiental. A comparação de terras privadas com terras públicas sem destinação não produz a mesma conclusão que sua comparação com áreas protegidas. Já \citeonline{ProbstEtAl2020}, ao examinar o programa Terra Legal, encontram aumento do desmatamento entre pequenos e médios proprietários titulados e maior resposta aos preços agropecuários. A regularização, desse modo, não deve ser confundida com um indicador inequívoco de conservação.

Esses resultados precisam ser lidos em função de seus objetos: Pacheco e Meyer comparam regimes fundiários em diferentes escalas e períodos, enquanto Probst e coautores examinam a resposta a uma iniciativa de titulação. A diferença entre uma comparação de regimes e a avaliação de uma intervenção impede tratar seus resultados como estimativas intercambiáveis. Sua contribuição para o presente artigo está em explicitar o condicionamento institucional dos incentivos. A investigação específica da apropriação irregular e da valorização especulativa da terra permanece fora de seu recorte.

A mesma cautela se aplica ao Cadastro Ambiental Rural. \citeonline{JungEtAl2022} analisam Pará e Mato Grosso e identificam respostas distintas da conversão para lavouras e pastagens, com heterogeneidade por porte dos estabelecimentos e conformidade ambiental. Em seu modelo dinâmico de uso da terra, com procedimentos para enfrentar a seleção no cadastro, não encontram efeito estatisticamente significativo da inscrição sobre a conversão para lavouras, enquanto identificam respostas da expansão de pastagens que variam entre os grupos examinados. Os resultados apontam o acesso ao crédito como um possível mecanismo. Essa evidência reforça a necessidade de examinar como instrumentos distintos afetam oportunidades produtivas e custos de conversão, em vez de reunir todo registro administrativo sob uma medida única de qualidade institucional.

O resultado não significativo para lavouras é distinto de ausência de estudos e não estabelece que o efeito seja exatamente zero. Tampouco permite concluir que o cadastro seja irrelevante para qualquer uso da terra. Em conjunto com os demais resultados do estudo, ele mostra por que o destino da conversão e as condições dos produtores importam para interpretar o alcance de uma intervenção.

Fora do contexto brasileiro, \citeonline{ViallonEtAl2019} mostram, em doze casos de uso da terra na Suíça, que a implementação local depende da capacidade das autoridades e dos direitos adquiridos pelos proprietários. \citeonline{AnderssonEtAl2014}, a partir de comunidades florestais na Bolívia, relacionam condições das florestas a diferentes combinações de elaboração de regras, monitoramento e sanções organizadas pelos usuários. Esses trabalhos ajudam a motivar a dimensão territorial da análise, não identificam, para Mato Grosso, a superioridade causal de um nível de governança sobre outro.

Isso não exclui, por óbvio, a relevância de instrumentos econômicos de conservação, como pagamentos por serviços ambientais, mercados de carbono ou financiamento ambiental. As receitas potenciais associadas à conservação do carbono florestal, por exemplo, podem modificar os retornos relativos enfrentados pelos agentes \cite{Angelsen2010,SticklerEtAl2013}. O recorte deste artigo, contudo, concentra-se nos preços agropecuários e nas condições institucionais que poderão ser observadas em séries territoriais comparáveis. Nenhuma conclusão sobre a eficácia de um instrumento específico será extraída sem observar sua aplicação e dispor de um desenho de avaliação apropriado.

Para o presente artigo, a contribuição institucional está, portanto, no exame da sensibilidade da conversão aos incentivos econômicos em contextos de implementação distintos, e não em classificar instituições como uniformemente favoráveis à conservação. Os resultados de \citeonline{ProbstEtAl2020} e \citeonline{JungEtAl2022} constituem antecedentes dessa análise. Em diálogo com eles, pretende-se distinguir as dimensões institucionais que os dados permitem observar e relacioná-las à resposta da conversão por destino agropecuário, considerando a capacidade produtiva local, sem equiparar registros administrativos à efetividade das regras.

% Origem: REF2, REF9 e comparação entre instrumentos em REF11.
\subsection{Fiscalização e pressão econômica em Mato Grosso}

A evidência disponível indica que monitoramento e fiscalização podem conter a conversão, embora a resposta aos incentivos econômicos varie conforme o instrumento e a área considerada. \citeonline{HardingEtAl2021} investigam três instrumentos na Amazônia brasileira, entre os quais a Moratória da Soja, que atua por critérios de compra estabelecidos pelas empresas da cadeia produtiva. Encontram redução do desmatamento associado aos preços nos municípios da lista prioritária, maior sensibilidade a preços de outros cultivos em municípios expostos à Moratória da Soja e ampliação da pressão sobre terras não protegidas associada às zonas de conservação. Os resultados sugerem observar também substituições e deslocamentos espaciais, em vez de limitar a avaliação ao efeito médio na área diretamente alcançada pela política. Essa diversidade de respostas fundamenta o exame do condicionamento institucional previsto em H1-b, sem pressupor que toda regulação reduza uniformemente a sensibilidade aos preços.

\citeonline{AssuncaoEtAl2023}, por sua vez, utilizam limitações de detecção por nuvens no sistema DETER como fonte de variação para identificar a atuação fiscalizatória. O estudo encontra efeito de contenção do desmatamento pelo monitoramento e pela fiscalização. Seu desenho evidencia por que a correlação entre número de autos e conversão é insuficiente: a fiscalização pode se concentrar justamente onde a pressão de desmatamento é maior. A estratégia com nuvens tampouco será importada automaticamente para este artigo, pois condições meteorológicas podem afetar a agricultura e a conversão por outras vias.

Em Mato Grosso, \citeonline{Sousa2016} examina multas, embargos e apreensões por meio de regressões e análise espacial, reportando correlações negativas com o desmatamento e diferenças de implementação entre sanções. O estudo destaca a baixa arrecadação das multas e a relevância dos embargos. Esses resultados constituem evidência sobre associações e funcionamento administrativo, sem dispensar o enfrentamento da seleção da fiscalização. Por sua vez, \citeonline{SticklerEtAl2013} analisam custos e consequências potenciais do cumprimento das obrigações de conservação previstas no Código Florestal no estado e discutem a instabilidade regulatória e a combinação com incentivos positivos à conservação.

Esses resultados não são propriamente incompatíveis: refletem desenhos voltados a diferentes margens da decisão de uso da terra, em períodos e ambientes institucionais distintos. Em diálogo com \citeonline{HardingEtAl2021}, a questão que permanece para este artigo é como observar, sob uma base territorial comum, se a sensibilidade da conversão aos preços varia com as condições de implementação das regras, a capacidade produtiva e o destino agropecuário da área convertida. Busca-se, assim, aproximar a análise da pressão econômica da heterogeneidade da fiscalização em Mato Grosso, sem reduzir essa última à contagem de autuações. Trata-se de delimitar uma contribuição empírica para Mato Grosso, sem afirmar ausência de estudos que já combinam preços e instituições.

% Origem: REF4--REF6, REF10 e contribuição/transição empírica de REF11.
\subsection{Heterogeneidade territorial e contribuição do artigo}

Em Mato Grosso, a expansão de duas safras anuais está associada ao acesso a transportes, à produção agrícola anterior e à conexão com mercados, como mostram \citeonline{VanWeyEtAl2013}. Os mapas de \citeonline{SimoesEtAl2020} permitem distinguir trajetórias de agricultura, pastagem e vegetação entre 2001 e 2017, abrangendo Amazônia, Cerrado e Pantanal. Essas contribuições ajudam a separar expansão territorial e intensificação: o aumento da produção ou da frequência de cultivo não informa, isoladamente, quanto de vegetação foi convertido. No Pantanal, \citeonline{SeidlEtAl2001} examinam indicadores econômicos e incentivos dos proprietários relacionados à expansão de pastagens cultivadas. Esse estudo abrange o Pantanal brasileiro e não deve ser tratado como estimativa específica dos municípios pantaneiros de Mato Grosso.

Tomadas em conjunto, essas contribuições justificam investigar a sensibilidade da conversão aos preços e seu condicionamento institucional com atenção às atividades, aos estoques remanescentes e aos biomas. As evidências revistas, contudo, empregam resultados de interesse distintos: perda florestal, conversão para usos específicos, mudanças de área e intensificação produtiva. Sua comparação requer preservar essas diferenças de mensuração. Para o presente artigo, a conversão bruta de vegetação nativa para agricultura e pastagem constitui o resultado de interesse, ao tempo em que mudanças líquidas de área não permitem recuperá-la quando coexistem fluxos de conversão, regeneração e substituição entre usos.

A dimensão intertemporal aparece no horizonte em que produtores percebem preços, incorrem em custos de conversão e antecipam a aplicação das regras. A conversão presente também modifica o estoque disponível para decisões posteriores. Isso dá conteúdo ao tratamento da vegetação como componente endógeno do processo econômico, sem confundir essa interpretação conceitual com a endogeneidade econométrica de uma variável explicativa. O horizonte das respostas deverá ser avaliado na etapa empírica, não se pressupõe que defasagens observadas revelem preferências temporais ou expectativas individuais.

Dito isso, a seção descritiva deverá caracterizar a distribuição e a evolução da vegetação nativa, da agricultura e das pastagens em Mato Grosso, distinguindo, quando as transições puderem ser construídas, a abertura de novas áreas da reorganização de usos já existentes. Deverá examinar também as diferenças territoriais de capacidade produtiva, a composição das atividades agropecuárias e a cobertura dos registros institucionais. Essas dimensões importam porque áreas com estoques remanescentes, possibilidades produtivas e condições de implementação distintas podem responder de maneira diferente aos mesmos incentivos. Sua comparação entre biomas dependerá da compatibilidade das classificações, da harmonização territorial e do suporte amostral.

Essa caracterização poderá motivar a investigação posterior das relações previstas em H1: como os incentivos econômicos se relacionam com a conversão; se a resposta aos preços varia conforme a capacidade produtiva local; e como o ambiente institucional condiciona essa relação. A descrição permitirá avaliar a variação disponível e os limites das medidas candidatas, mas não demonstrará identificação causal. A unidade município--ano e o período 2001--2024 permanecem candidatos, condicionados à construção e à compatibilização dos dados. A escolha das especificações e das condições de interpretação caberá à seção de estratégia econométrica, após esse exame.

Assim, o interesse de Mato Grosso para a pesquisa não decorre apenas da presença de três biomas, mas da possibilidade de comparar respostas à pressão econômica sob diferentes combinações de capacidade produtiva, usos agropecuários e condições institucionais. Essa comparação procura esclarecer o alcance territorial das relações previstas em H1, em vez de assumir que resultados encontrados em uma atividade ou contexto se aplicam aos demais. Seu valor analítico dependerá das diferenças que puderem ser estimadas com precisão, e não da diversidade geográfica tomada isoladamente.

Em síntese, busca-se avançar na compreensão de como condições econômicas e institucionais se relacionam com a formação e a transformação do estoque ambiental, preservando a distinção entre mecanismos conceituais, associações observadas e efeitos cuja interpretação causal possa ser sustentada. A contribuição pretendida resulta da articulação entre a mensuração da conversão bruta por destino, a heterogeneidade da resposta aos preços conforme a capacidade produtiva e seu condicionamento institucional. Essa articulação poderá informar o debate sobre a adequação territorial das políticas de conservação, sem pressupor que a reunião dessas dimensões constitua, por si só, inovação metodológica ou identificação causal. Resultados nulos ou contrários às previsões também deverão integrar essa avaliação, considerando sua precisão, o contexto em que foram obtidos e a correspondência entre as medidas utilizadas e as hipóteses do artigo.
